\documentclass[11pt,a4paper]{article}
\usepackage[T1]{fontenc}
\usepackage[utf8]{inputenc}
\usepackage{amsmath,amssymb,amsthm}
\usepackage[margin=1in]{geometry}
\usepackage[colorlinks=true,linkcolor=blue,urlcolor=blue]{hyperref}
\theoremstyle{plain}
\newtheorem{theorem}{Theorem}
\newtheorem{lemma}[theorem]{Lemma}
\newtheorem{proposition}[theorem]{Proposition}
\newtheorem{corollary}[theorem]{Corollary}
\theoremstyle{definition}
\newtheorem{remark}[theorem]{Remark}

\title{A cross-alphabet identity for bounded-difference square arrays: OEIS A125536}
\author{Adrian Perez Fontelles\\ \small Independent researcher}
\date{14 September 2026}
\begin{document}
\maketitle

\begin{abstract}
OEIS A125536 carries an empirical identity relating its counts to those of the same problem over an
alphabet one smaller. It is true, and sharply so: the threshold the entry states is exactly the
diameter bound of the grid graph, and the identity fails below it. The proof is four lines and
uses no machinery.
\end{abstract}

\noindent\small 2020 Mathematics Subject Classification. 05A15, 05B45.\normalsize

\section{The conjecture}

OEIS A125536 is ``Number of n X n arrays with entries in 1..6 in which adjacent entries differ by 4 or less (adjacent means in x or y directions).''. It has offset $0$ and begins
\[
1, 6, 1058, 5748554, 948155870210, 4756811244059748042, 725704975038742796696062050, 3366853797682513115652233385161124618,\ \dots
\]
The entry states, as a conjecture contributed by R. H. Hardin and never marked settled:
\begin{quote}
[Empirical] a(base+1,n,diff)=a(base,n,diff)+F(n,diff) for base>=2.diff.(n-1) and F(n,diff)= (n=2, A063496(diff+1)) (n=3, A068744(diff+1)) (n=4, A068745(diff+1)) (n=5, A068746(diff+1)) (n=6, A068747(diff+1))
\end{quote}
As of the ``Last modified'' line on the live entry (Mar 31 2012 12:35:05, revision 6)
this is still recorded as empirical. For this entry's own alphabet the threshold $2d(n-1)\le b$ leaves only $n=1$, so the identity is not a statement about this entry's published terms beyond the first; it is a statement about the family, and it is the family statement that is proved below.

\section{Notation}

For integers $b\ge1$, $d\ge1$ and $n\ge1$ write $T(b,d,n)$ for the number of $n\times n$ arrays
with entries in $\{1,\dots,b\}$ in which any two orthogonally adjacent entries differ by at
most $d$. The entry is $T(6,4,n)$, and the conjecture is
\[
T(b+1,d,n)\;=\;T(b,d,n)+F(n,d)\qquad\text{for } b\ \ge\ 2d(n-1),
\]
with $F(n,d)$ the quantity the entry tabulates through its reference sequences ($n=2$: OEIS A063496, $n=3$: OEIS A068744, $n=4$: OEIS A068745, $n=5$: OEIS A068746, $n=6$: OEIS A068747).

\section{The proof}

\begin{theorem}
For all $n\ge1$, $d\ge1$ and $b\ge 2d(n-1)$,
\[
T(b+1,d,n)-T(b,d,n)\;=\;F(n,d),
\]
where $F(n,d)$ is the number of admissible arrays over $\mathbb{Z}$ whose maximum entry is
$0$ --- a quantity depending on $n$ and $d$ only.
\end{theorem}

\begin{proof}
\emph{Step 1.} An admissible array over $\{1,\dots,b+1\}$ either uses the value $b+1$ or
does not, and those that do not are exactly the admissible arrays over $\{1,\dots,b\}$. So
$T(b+1,d,n)-T(b,d,n)$ counts the admissible arrays over $\{1,\dots,b+1\}$ that use $b+1$.

\emph{Step 2.} Every admissible array satisfies $\max-\min\le 2d(n-1)$. Let $u$ attain the
maximum and $v$ the minimum. In the grid graph on $n\times n$ cells with orthogonal adjacency
any two cells are joined by a path of length at most $2(n-1)$ --- move along a row, then along
a column --- and along a path of length $L$ the entry changes by at most $dL$, since each step
is between adjacent cells.

\emph{Step 3.} An array counted in Step 1 has maximum exactly $b+1$, so by Step 2 all of its
entries lie in $[\,b+1-2d(n-1),\ b+1\,]$. When $b\ge 2d(n-1)$ that interval is contained in
$\{1,\dots,b+1\}$, so the constraint ``entries in $\{1,\dots,b+1\}$'' is implied by the
maximum and the range bound and imposes nothing further. The arrays counted in Step 1 are
therefore in bijection --- by the translation carrying the maximum to $0$ --- with the
admissible arrays over $\mathbb{Z}$ of maximum $0$, and that count involves no $b$.
\end{proof}

\begin{remark}
The threshold is sharp, and not merely sufficient. Below it the window no longer fits and the
identity fails: at $b=5$, $d=3$, $n=3$ the difference $T(6,3,3)-T(5,3,3)$ is $964{,}755$
while $F(3,3)=1{,}253{,}329$.
\end{remark}

\section{Verification}

Four checks, each independent of the algebra above.

First, the count $T(6,4,n)$ was computed by an independent row-pair dynamic program and
reproduces the $4$ published terms of A125536 that the computation reaches,
exactly, in integer arithmetic.

Second, the identity was evaluated directly at the entry's own alphabet size for every $n$ the
threshold covers, and at one $n$ beyond it, where it fails as the remark says.

Third, $F(n,d)$ as computed here agrees with the entry's own reference sequences. The entry
writes them as ``$\mathrm{A}\dots(\mathit{diff}+1)$'', and the index is positional --- the
$(\mathit{diff}+1)$st term as listed. That reading is forced: the sequences do not share an
offset, and only the positional reading makes $n=2$ and $n=3$ agree simultaneously.

Fourth, the following table was produced by the same program and can be checked by hand at
small $n$:
\[
\begin{array}{rrrr}
n & T(6+1,4,n) & T(6,4,n) & F(n,4) \\ \hline
1 & 7 & 6 & 1 \\
2 & 1527 & 1058 & 489 \\
\end{array}
\]

\begin{thebibliography}{9}
\bibitem{oeis} The OEIS Foundation, \emph{The On-Line Encyclopedia of Integer Sequences},
\url{https://oeis.org}, 2026. Sequence A125536: \url{https://oeis.org/A125536}.
\end{thebibliography}

\end{document}
